\section{Calibration and Estimation}
\label{sec:calibration}

This section describes how the calibrated notch model of Section~\ref{sec:model}
is taken to the data. I proceed in three steps: estimate the curvature
parameter $\alpha$ of the cost specification from firms away from the notch;
assign each firm a cost/productivity parameter $A_i$ as a calibration input; and
extract the headline behavioural elasticity, reporting bootstrap standard errors
for the bunching estimates.

\subsection{Estimating the cost-function curvature}

The cost specification is Cobb--Douglas, $Y = A\,X^{\alpha}$, so that in logs
\(\ln Y = \ln A + \alpha \ln X\). I estimate $\alpha$ by a weighted regression
of $\ln Y$ on $\ln X$ across firms. To avoid contaminating the estimate with the
behavioural response the model is meant to explain, the estimation sample excludes
firms in the manipulation region around the data-year threshold
$T^{*}=\pounds85{,}000$, retaining firms with turnover below \pounds75{,}000 or
above \pounds95{,}000, which face no marginal registration incentive.

The estimated curvature is $\alpha \approx 0.989$ on the \pounds85{,}000 synthetic
data (precisely $0.9888$; an earlier hardcoded value of $0.9864$ is no longer
used)---near-constant returns.\footnote{This is an un-instrumented OLS slope of
$\ln Y$ on $\ln X$ with a single input, with no control-function or
proxy-variable correction (Levinsohn--Petrin, Olley--Pakes, ACF). If more
productive firms purchase more inputs, the estimate is biased upward towards
unity, so any object that turns on the gap between $\alpha$ and $1$ is fragile.
I therefore route the load-bearing objects away from $\alpha$: the dominated
region is the exact Kleven--Waseem width, and the marginal buncher is evaluated in
the iso-elastic representation with the behavioural elasticity
(Section~\ref{ssec:buncher}).}

\subsection{Assigning the firm-level cost parameter}

Given $\alpha$, I assign each firm the cost/productivity parameter
$A_i = Y_i / X_i^{\alpha}$, the production residual implied by its observed
turnover $Y_i$ and inputs $X_i$. This is a \emph{calibration input} for the
forward re-solution of Section~\ref{sec:model}, not a claim of structural
identification: it is whatever cost level rationalises the firm's observed
allocation at the \pounds85{,}000 schedule, held fixed when the schedule changes.
I do not pre-process turnover with any ``unbunching'' rule or inject uncertainty
before forming $A_i$, so the same residual is carried unchanged into the forward
solution. The assignment is available for the $91.5\%$ of firms with strictly
positive observed inputs and turnover; the values are heterogeneous and broadly
increasing in firm size.

\subsection{The headline behavioural elasticity}
\label{ssec:elasticities}

The behavioural object of interest is the elasticity of reported turnover with
respect to the net-of-tax rate. I obtain it by perturbing the VAT rate by a
small $\Delta\tau$, re-solving each firm's problem, and computing
\[
  \varepsilon
  = \frac{\ln y^{*}(\tau_{\max}+\Delta\tau) - \ln y^{*}(\tau_{\max})}
         {\ln(1-\tau_{\max}-\Delta\tau) - \ln(1-\tau_{\max})} .
\]
Computed firm by firm near the threshold, this local turnover elasticity has a
median of approximately $\mathbf{0.17}$ and a weighted mean of approximately
$0.32$ (the mean pulled up by a thin tail of very responsive firms). As a loose
magnitude check, the median is of the same order as the $0.09$--$0.18$
VAT-\emph{base} elasticity reported by \citet{liuetal2021} for the UK; this is
only a magnitude comparison, since their object is a value-added (base) elasticity
for which they decline to back out a turnover elasticity, whereas mine is a
turnover elasticity. This elasticity governs the firm-level re-optimisation behind
the reform costs of Section~\ref{ssec:schedule_costs} and gives the behavioural
reading of the bunching evidence; a fuller behavioural counterfactual that
headlines re-optimised magnitudes I leave to future work. Alongside it I report the
descriptive bunching ratio $b=0.060$
(SE $0.005$)---the excess-mass \emph{width} $\Delta y^{*}/y^{*}$, not an
elasticity.\footnote{The bunching estimator also implies a substitution
(CES-wedge) elasticity $\sigma\approx4.85$ (SE $0.12$), which reads the full
excess span as a response; it is not comparable to the local turnover elasticity.
See Appendix~\ref{app:inference}.}

The mass-conservation constraint (Appendix~\ref{app:inference}) raises $b$ (from
$0.054$) and the excess mass $E$ (from $8{,}179$ to $8{,}712$, SE $662$) relative
to the unconstrained fit, and forces excess and missing mass into agreement
($E=8{,}712$ versus $\Delta_{R}=9{,}938$, against $8{,}179$ versus $17{,}792$
unconstrained); the implied displaced share $\Pi=0.211$ (SE $0.005$) is robust
even where the wedge level is not. The sigmoid steepness $k$ does not affect these
data-side magnitudes but governs the width of the structural schedule's dominated
region; I set $k=0.001$, whose implied band ($\pounds19.7$k) is closest to the
exact \citet{klevenwaseem2013} notch region $T^{*}\tau/(1-\tau)=\pounds21.25$k.
Re-solving the calibrated model reproduces the observed near-threshold
distribution closely, but this is true essentially by construction (the model is
fitted to that distribution) and I do not treat the fit as an independent result.
